\begin{definition}[CSP {template/instance/assignment}] {We define a \emph{CSP template} to be a finite set of predicates $\Gamma =\{P_1, P_2, \ldots, P_l\}$ over a finite domain $D$, where for all $i \in [\ell]$, $P_i$ has arity $k_i \ge 0$.}
{An \emph{instance} $\Phi=(V,\mathcal{C})$ of the constraint satisfaction problem (CSP) associated with the template $\Gamma$ (i.e., ``$\Phi$ is an instance of $\CSP(\Gamma)$''), consists of a variable set $V$ and a constraint set $\mathcal C$. The variable set consists of} $n$ variables $V =\{u_1, u_2, \ldots, u_n\}.$ {The} $m$ constraints $\mathcal{C}=\{C_1, C_2, \ldots, C_m\}$ 
each {consist} of a tuple of variables $C_j = (u_{j,1}, u_{j,2}, \ldots, u_{j,{l_j}}) \in V^{l_j}$ and an associated predicate $P^{(j)} \in \Gamma$ of the same arity $l_j$. 
An assignment $\sigma : V \rightarrow D$ is said to satisfy the constraint $C_j$ if $\sigma(C_j)=(\sigma(u_{j,1}), \sigma(u_{j,2}), \ldots, \sigma(u_{j,{l_j}})) \in P^{(j)}$.
\end{definition}

{We now define two computational problems associated with $\CSP(\Gamma)$ where an instance $\Phi$ of $\CSP(\Gamma)$ is given as input.

\begin{definition}[{CSP decision/search problems}] In the \emph{decision version} of $\CSP(\Gamma)$, the objective is to decide if there is an assignment to a given instance of $\CSP(\Gamma)$ that satisfies all the constraints. In the \emph{search version} of $\CSP(\Gamma)$, the objective is to output an assignment to a given instance of $\CSP(\Gamma)$ that satisfies all the constraints.
\end{definition}}

We next define {promise constraint satisfaction problems} (PCSP).  {We begin by defining what a promise template is.}\footnote{{Other works, such as \cite{BBKO21}, parameterize the promise CSP as a pair of a relational structures with a homomorphism between them.}}

{
\begin{definition}[promise template]
    Let $D_1, D_2$ be finite domains. A promise template consists of a finite collection $\Gamma = \{ (P_1,Q_1), (P_2,Q_2),\ldots, (P_l,Q_l)\}$ of pairs of predicates, where for all $i \in [\ell]$, $P_i$ has arity $k_i\ge 0$ over domain $D_1$ and each $Q_i$ has arity $k_i$ over domain $D_2$. Furthermore, we assert there exists a map $h:D_1\rightarrow D_2$ such that for all $i\in[l]$ and $\textbf{x} \in D_1^{k_i}$, $\textbf{x} \in P_i$ implies $h(\textbf{x}) \in Q_i$. In other words, $h$ is a homomorphism from $P_i$ to $Q_i$ for all $i \in [l]$.
\end{definition}
}

We say that $(D_1,D_2)$ is the domain of the promise template $\Gamma$.

\begin{definition}[{PCSP instance/assignments}]
	In an instance $\Phi=(V,\mathcal{C})$ of $\PCSP(\Gamma)$, we have a set of $n$ variables $V =\{u_1, u_2, \ldots, u_n\}$ and $m$ constraints $\mathcal{C}=\{C_1, C_2, \ldots, C_m\}$ 
each consisting of a tuple of variables $C_j = (u_{j,1}, u_{j,2}, \ldots, u_{j,{l_j}}) \in V^{l_j}$ and an associated predicate pair $(P^{(j)},Q^{(j)}) \in \Gamma$ of the same arity $l_j$.

{A $P$-assignment is a map $\sigma_1 : V \rightarrow D_1$ and a $Q$-assignment is a map $\sigma_2 : V \rightarrow D_2$.} A {$P$}-assignment {$\sigma_1$} is said to strongly satisfy the constraint $C_j$ if $\sigma(C_j)\in P^{(j)}$, and {a $Q$-assignment $\sigma_2$} is said to weakly satisfy the constraint $C_j$ if $\sigma(C_j) \in Q^{(j)}$.
\end{definition}

{Crucially, due to the existence of the homomorphism $h$ connecting each $P_i$ with its corresponding $Q_i$, any $P$-assignment strongly satisfying a collection of constraint can be converted into a $Q$-assignment weakly satisfying those constraints by composing $h$ with the $P$-assignment.} {If $\Gamma = \{(P,Q)\}$ consists a single pair of predicates, we often state $\Gamma = (P,Q)$ and write $\PCSP(P, Q)$ for $\PCSP(\Gamma)$.} The following are computational problems associated with {instances of} $\PCSP(\Gamma)$. 

\begin{definition}[{PCSP problems}]\label{def:pcsp-prob} In the decision version of $\PCSP(\Gamma)$, given an input instance $\Phi=(V,\mathcal{C})$ of $\PCSP(\Gamma)$, the objective is to distinguish between the two cases.
  \begin{enumerate}
\item  There is a $P$-assignment $\sigma_1 : V \rightarrow D_1$ that strongly satisfies all the constraints.
\item There is no $Q$-assignment $\sigma_2 : V \rightarrow D_2$ that weakly satisfies all the constraints.  
\end{enumerate}

In the search version of $\PCSP(\Gamma)$, given an input instance $\Phi=(V,\mathcal{C})$ of $\PCSP(\Gamma)$ with the promise that there is an assignment $\sigma_1: V\rightarrow D_1$ that strongly satisfies all the constraints, the objective is to find an assignment $\sigma_2: V \rightarrow D_2$ that weakly satisfies all the constraints.\footnote{{This explains the terminology \emph{promise} CSPs to refer to such problems. For the decision version, a promise problem consists of disjoint sets of Yes and No instances, which together need not exhaust all instances.}}
 
{An algorithm for $\PCSP(\Gamma)$ is said to be $f$-robust algorithm for a function $f:\R \rightarrow \R$ if the following holds. The inputs to the algorithm are an instance $\Phi=(V,\mathcal{C})$ of $\PCSP(\Gamma)$ and a parameter $\epsilon$} with the promise that there is an assignment $\sigma_1: V\rightarrow D_1$ that strongly satisfies $1-\epsilon$ fraction of the constraints. {The objective of the algorithm is to output an assignment} $\sigma_2: V \rightarrow D_2$ that weakly satisfies at least $1-f(\epsilon)$ fraction of the constraints. {We call an algorithm a robust algorithm for $\PCSP(\Gamma)$ if it is $f$-robust for some monotone and non-negative function $f$ with $f(\epsilon)\rightarrow 0$ as $\epsilon \rightarrow 0$.}
\end{definition}

In this paper, we restrict ourselves to PCSP \emph{templates} where both the domains are equal to the set ${\B :=} \{-1,+1\}$. {We call such templates \emph{Boolean} templates.} {For technical reasons, to more closely align with the conventions in the robust CSP literature, we consider a different notion of an PCSP instance which incorporates \emph{variable folding}, i.e., the negation of variables. We also assume that the PCSP instance allows for the setting of constants, which we call \emph{idempotence}.}

\begin{definition}[Boolean folded idempotent {instance/assignment}] {Given a Boolean template $\Gamma=\{(P_1,Q_1),$ $(P_2,Q_2),\ldots,$ $(P_l,Q_l)\}$ where $P_i \subseteq Q_i \subseteq \B^{k_i}$ for every $i \in [l]$, an instance $\Phi=(V,\mathcal{C})$ of $\fiPCSP(\Gamma)$ consists of a set of $n$ variables $V = \{u_1, \hdots, u_n\}$ and a set of $m$ constraints $\mathcal{C} = \{C_1, \hdots, C_m\}$. For $i \in [n]$, each variable $u_i$ has two corresponding \emph{literals}: $u_i$ and $\overline{u_i}$, the positive and negative literals, respectively. {Observe that we identify the positive literal with the variable itself.} We also have two constant literals corresponding to the elements of $\B$. For each $j \in [m]$, $C_j$ consists} of a tuple of literals $C_j = (x_{j,1}, x_{j,2}, \ldots, x_{j,{l_j}}) \in {(}V\cup \overline{V} \cup \B{)}^{l_j}$ and an associated predicate pair $(P^{(j)},Q^{(j)}) \in \Gamma$ of the same arity $l_j$. 

Given an assignment\footnote{{Since the families $P_i$ and $Q_i$ of predicates have a common domain $\B$, we need not distinguish between ``$P$-assignments'' and ``$Q$-assignments'' in the Boolean setting.}} $\sigma : V \rightarrow \B,$ {define its extended assignment $\sigma' : V \cup \overline{V} \cup \B \rightarrow \B$ to be} $\sigma'(u_i)=\sigma(u_i)$ and $\sigma'(\overline{u_i})=-\sigma(u_i)$ for every $i \in [n]$ {and $\sigma'(b) = b$ for $b \in \B$}. The assignment $\sigma$ is said to strongly (and {respectively} weakly) satisfy the constraint $C_j$ in $\Phi$ if $\sigma'(C_j) \in P^{(j)}$ (and {respectively} $\sigma'(C_j) \in Q^{(j)}$). 
\end{definition}
